% Generated from the Lean dependency data; see the comment at the top of content.tex.
\chapter{The multiplicity formula}
\label{chap:multformula}
Thesis Ch.~7. The multiplicity formula is proved by strong induction on $\dim G$, together with the
geometric expansion of the distribution $J$. Test functions are taken in $C_c^\infty$ modulo $A_G(F)$
instead of the Harish-Chandra Schwartz space; every stated theorem of the trace formula uses this space.

\section{\texorpdfstring{Strongly cuspidal functions and the distribution $J$}{Strongly cuspidal functions and the distribution J}}
\begin{definition}[Strongly cuspidal test functions]
\label{ScuspTestFn}
\lean{GSpinGGP.ScuspTestFn}
\leanok
\uses{CentralCharOn, G, GGPData, GSpin, H, N}
$\scusp(A_G(F)\backslash G(F),\chi^{-1})$, $C_c^\infty$ version: locally constant $f$ on $G(F)$ with
$f(ag)=\chi(a)^{-1}f(g)$ for $a\in A_G(F)$, compactly supported modulo $A_G(F)$, and strongly cuspidal
($\int_{U(F)}f(pu)\,du=0$ for every proper parabolic $P=MU$ and $p\in P(F)$) (Thesis Ch.~1, \S4.5, \S7.1).
\end{definition}

\begin{definition}[The distribution $J$]
\label{distribJ}
\lean{GSpinGGP.distribJ}
\leanok
\uses{G, GGPData, GSpin, H, N, xi}
\[ J(f)=\int_{H(F)\backslash G(F)/A_G(F)}K(f,x)\,dx,\qquad K(f,x)=\int_{H(F)/A_{G,H}(F)}f(x^{-1}hx)\xi(h)\,dh \]
(Thesis \S7.1, \S10.1), for the Weil quotient measure on $H(F)\backslash G(F)/A_G(F)$.
\end{definition}

\begin{theorem}[Convergence of $J$]
\label{distribJ_convergent}
\lean{GSpinGGP.GGPData.distribJ_convergent}
\leanok
\uses{CentralCharOn, G, GGPData, GSpin, H, N, ScuspTestFn, xi}
For $f\in\scusp(A_G\backslash G,\chi^{-1})$, $x\mapsto K(f,x)$ is integrable on $H(F)\backslash G(F)/A_G(F)$
(Thesis Thm.~10.1.1(ii) with Prop.~8.4.1(iii); = Thm.~1.0.3).
\end{theorem}
\begin{proof}
Not yet formalized. The bound $|K(f,x)|\le\nu(f)\,\Xi^{H\backslash G}(x)^2\sigma(x)^{-d}$ for strongly cuspidal $f$,
and integrability of $\Xi^{H\backslash G}(x)^2\sigma(x)^{-d}$ for $d$ large.
\end{proof}

\begin{lemma}[$J$ is linear]
\label{distribJ_add}
\lean{GSpinGGP.GGPData.distribJ_add}
\leanok
\uses{CentralCharOn, G, GGPData, ScuspTestFn, distribJ}
$J(f+g)=J(f)+J(g)$ for $f,g\in\scusp(A_G\backslash G,\chi^{-1})$ (Thesis \S10.1).
\end{lemma}
\begin{proof}
\leanok
\uses{GSpin, H, N, distribJ_convergent, xi}
Linearity of the integral, given convergence.
\end{proof}

\begin{definition}[The quasi-character $\theta_f$]
\label{thetaF}
\lean{GSpinGGP.thetaF}
\leanok
\uses{G, GGPData, GSpin, H, N, gspin}
For strongly cuspidal $f$, $\theta_f$ is the function representing the orbital pairing
$f'\mapsto\int_{G(F)/A_G(F)}\int_{G(F)}f(x^{-1}gx)f'(g)\,dg\,dx$ (Thesis \S4.8, \cite{BP} \S5.5);
$0$ if there is none. Convention: $\theta_f(ag)=\chi(a)^{-1}\theta_f(g)$.
\end{definition}

\begin{literature}[$\theta_f$ represents the orbital pairing]
\label{LIT-bp_thetaF_exists}
\lean{GSpinGGP.LitTF.bp_thetaF_exists}
\leanok
\uses{G, GGPData, GSpin, IsQuasiCharacter, gspin}
Let $f_0\in C_c^\infty(G(F))$ be strongly cuspidal. Then there is a quasi-character $\theta$, locally integrable
and locally constant on $\Greg(F)$, with $\int_{G(F)}\theta f'\,dg=\int_{G(F)/A_G(F)}\int_{G(F)}f_0(x^{-1}gx)f'(g)\,dg\,dx$
for every $f'\in C_c^\infty(G(F))$, the outer integrand being integrable.

\assumedfrom{\cite{BP}, Thm.~5.5.1(ii),(iv), Prop.~5.6.1(ii) (= Thesis Thm.~4.7.1, Prop.~4.8.1(ii)); \cite{Art91}; Weyl integration formula}
\end{literature}

\begin{lemma}[$\chi$-transfer of $\theta_f$]
\label{isThetaF_chiProj}
\lean{GSpinGGP.isThetaF_chiProj}
\leanok
\uses{CentralCharOn, G, GGPData, GSpin, H, IsQuasiCharacter, N, gspin}
For strongly cuspidal $f_0\in C_c^\infty(G(F))$, the orbital pairing of its $\chi$-projection $P_\chi f_0$
is represented by a quasi-character (Thesis \S4.8, \S7.1).
\end{lemma}
\begin{proof}
\leanok
\uses{LIT-bp_thetaF_exists, ScuspTestFn, jHat}
Glue central translates of the quasi-character of a cut-off of $P_\chi f_0$ along the fibres of the Harish-Chandra map $H_G$.
\end{proof}

\begin{proposition}[$\theta_f$ of a strongly cuspidal function]
\label{thetaF_spec}
\lean{GSpinGGP.thetaF_spec}
\leanok
\uses{CentralCharOn, G, GGPData, GSpin, H, IsQuasiCharacter, N, ScuspTestFn, gspin, thetaF}
For $f\in\scusp(A_G\backslash G,\chi^{-1})$, $\theta_f$ represents the orbital pairing of $f$ and is a
quasi-character (Thesis \S4.8, Prop.~4.8.1(ii)).
\end{proposition}
\begin{proof}
\leanok
\uses{LIT-hc_regularSet_open_conull, isThetaF_chiProj, jHat, weylDisc}
From the $\chi$-transfer and uniqueness on the open conull regular set.
\end{proof}

\begin{definition}[$\QC(G(F),\chi^{-1})$]
\label{QC}
\lean{GSpinGGP.QC}
\leanok
\uses{CentralCharOn, G, GGPData, GSpin, IsQuasiCharacter, gspin, jHat}
Quasi-characters $\theta$ on $G(F)$ with $\theta(ag)=\chi(a)^{-1}\theta(g)$ for $a\in A_G(F)$, $g\in\Greg(F)$
(Thesis \S4.1.1, \S7.2).
\end{definition}

\begin{lemma}[$\theta_f\in\QC(G,\chi^{-1})$]
\label{thetaF_mem_QC}
\lean{GSpinGGP.thetaF_mem_QC}
\leanok
\uses{CentralCharOn, G, GGPData, QC, ScuspTestFn, thetaF}
For $f\in\scusp(A_G\backslash G,\chi^{-1})$, $\theta_f\in\QC(G(F),\chi^{-1})$ (Thesis Prop.~4.8.1(ii)).
\end{lemma}
\begin{proof}
\leanok
\uses{GSpin, H, IsQuasiCharacter, LIT-hc_regularSet_open_conull, N, gspin, thetaF_spec}
From the previous proposition and central equivariance of $f$.
\end{proof}

\begin{lemma}[Characters are quasi-characters]
\label{charG_mem_QC}
\lean{GSpinGGP.charG_mem_QC}
\leanok
\uses{CentralCharOn, G, GGPData, GSpin, IsTempered, QC, UnitaryRep, charG, mult, orbitChoice}
If $\pi$ is irreducible tempered with central character $\eta$ on $A_G(F)$ then $\theta_\pi\in\QC(G(F),\eta)$
(Thesis Ch.~1, \S4.4).
\end{lemma}
\begin{proof}
\leanok
\uses{IsQuasiCharacter, LIT-hc_character_exists, LIT-hc_character_isQuasiCharacter, LIT-hc_regularSet_open_conull, gspin}
Harish-Chandra's local expansion, and central translation of characters.
\end{proof}

\section{The induction}
Thesis \S7.2 and \S7.4.

\begin{definition}[The geometric expansion at $(G,\chi)$]
\label{GeomAt}
\lean{GSpinGGP.GeomAt}
\leanok
\uses{CentralCharOn, G, GGPData, ScuspTestFn, distribJ, mGeom, thetaF}
$J(f)=\mgeom(\theta_f)$ for every $f\in\scusp(A_G\backslash G,\chi^{-1})$ (Thesis Thm.~7.2.1 = Thm.~12.2.1 = Thm.~1.0.5).
\end{definition}

\begin{definition}[The multiplicity formula at one triple]
\label{MultAt}
\lean{GSpinGGP.MultAt}
\leanok
\uses{G, GGPData, GSpin, IsTempered, UnitaryRep, charG, mGeom, mult}
For every irreducible tempered $\pi$ of $G(F)$ trivial on $A_{G,H}(F)$: $m(\pi)<\infty$ and
$m(\pi)=\mgeom(\theta_\pi)$.
\end{definition}

\begin{definition}[The induction hypothesis (HYP)]
\label{HYP}
\lean{GSpinGGP.HYP}
\leanok
\uses{CentralCharOn, GGPData, GeomAt, MultAt}
The geometric expansion (for every central character) and the multiplicity formula hold for every GGP
triple $(G',H',\xi')$ with $\dim G'<\dim G$ (Thesis \S7.2).
\end{definition}

\begin{definition}[Virtual tempered representations]
\label{VirtualTemp}
\lean{GSpinGGP.VirtualTemp}
\leanok
\uses{CentralCharOn, G, GGPData, GSpin, IsTempered, UnitaryRep, charG, mult, orbitChoice}
$\Rtemp(G,\eta)\otimes\mathbb C$: finite formal linear combinations of classes of irreducible tempered
representations with central character $\eta$ on $A_G(F)$ (Thesis \S3.4, \S7.4), with the linear extensions
$\theta_\Pi$, $m(\Pi)$, $\mgeom(\Pi)$.
\end{definition}

\begin{definition}[The defect]
\label{defect}
\lean{GSpinGGP.defect}
\leanok
\uses{CentralCharOn, G, GGPData, GSpin, IsTempered, UnitaryRep, VirtualTemp, charG, mGeom, mult, orbitChoice}
$\Pi\mapsto m(\Pi)-\mgeom(\Pi)$, extended linearly (Thesis Prop.~7.4.1(ii)). The multiplicity formula
for $\pi$ is $\mathrm{defect}[\pi]=0$.
\end{definition}

\begin{literature}[Elliptic isolation]
\label{LIT-bp_scusp_surj_ell}
\lean{GSpinGGP.LitTF.bp_scusp_surj_ell}
\leanok
\uses{G, GGPData, GSpin, IsQuasiCharacter, gspin}
There is an open, closed, invariant, $A_G(F)$-stable $\Omega\subset G(F)$ containing the elliptic regular set
and compact modulo conjugation and $A_G(F)$, such that every quasi-character supported (compactly
modulo conjugation) in $\Omega$ is $\theta_{f_0}$ on $\Greg(F)$ for some strongly cuspidal $f_0\in C_c^\infty(G(F))$
supported in $\Omega$.

\assumedfrom{\cite{BP}, Cor.~5.7.2(i) (= Thesis Cor.~4.8.3(i)), stated there for compact centre; adapted to a nontrivial split centre $A_G$ by its local proof}
\end{literature}

\begin{proposition}[Elliptic isolation, $\chi$-equivariant]
\label{exists_scusp_ell}
\lean{GSpinGGP.exists_scusp_ell}
\leanok
\uses{CentralCharOn, G, GGPData, GSpin, IsTempered, ScuspTestFn, UnitaryRep, VirtualTemp, charG, gspin, mult, orbitChoice, thetaF}
For every virtual tempered $\Pi$ with central character $\eta$ there is $f\in\scusp(A_G\backslash G,\eta)$ with
$\theta_f=\theta_\Pi$ on $\ellreg(F)$ (Thesis Cor.~4.8.3(i) in $\chi$-equivariant form).
This differs from the Thesis, which uses \cite{BP} Cor.~5.7.2(iv) and Schwartz functions.
\end{proposition}
\begin{proof}
\leanok
\uses{H, IsQuasiCharacter, LIT-bp_scusp_surj_ell, LIT-bp_thetaF_exists, LIT-hc_regularSet_open_conull, N, QC, charG_mem_QC, jHat}
Cut $\theta_\Pi$ off near the elliptic set, apply elliptic isolation, and project with $P_\chi$.
\end{proof}

\begin{proposition}[Thesis Prop.~7.4.1(ii), elliptic form]
\label{prop_11_5_1_ii_ell}
\lean{GSpinGGP.prop_11_5_1_ii_ell}
\leanok
\uses{CentralCharOn, G, GGPData, GSpin, HYP, IsTempered, ScuspTestFn, UnitaryRep, VirtualTemp, charG, defect, distribJ, gspin, mGeom, mult, orbitChoice, thetaF}
Assume (HYP). Let $f\in\scusp(A_G\backslash G,\eta)$ and $\Pi$ virtual tempered with central character $\eta$ such
that $\theta_\Pi=\theta_f$ on $\ellreg(F)$. Then $J(f)-\mgeom(\theta_f)=m(\Pi)-\mgeom(\Pi)$.
\end{proposition}
\begin{proof}
\leanok
\uses{defect_eq_zero_of_vChar_eq_zero_ell, exists_virtual_spectral_ell}
The spectral expansion gives a $\Pi_f$ with the same elliptic character and
$J(f)-\mgeom(\theta_f)=\mathrm{defect}(\Pi_f)$; $\Pi-\Pi_f$ has vanishing elliptic character, so it lies in $\Rind$,
where the defect vanishes (Chapter~\ref{chap:spectral}).
\end{proof}

\begin{proposition}[Geometric expansion implies multiplicity formula]
\label{multAt_of_hyp_geom}
\lean{GSpinGGP.multAt_of_hyp_geom}
\leanok
\uses{CentralCharOn, GGPData, GeomAt, HYP, MultAt}
(Thesis Prop.~7.4.1(iii), direction Thm.~7.2.1 $\Rightarrow$ Thm.~7.2.2.) Under (HYP), if the geometric
expansion holds at $G$ for all central characters, the multiplicity formula holds at $G$.
\end{proposition}
\begin{proof}
\leanok
\uses{G, GSpin, H, IsTempered, LIT-emoryTakeda_mult_le_one, LIT-schur_central, ScuspTestFn, UnitaryRep, VirtualTemp, charG, defect, distribJ, exists_scusp_ell, gspin, mGeom, mult, orbitChoice, prop_11_5_1_ii_ell, thetaF, xi}
Let $\eta$ be the central character of $\pi$ (Schur) and choose $f$ with $\theta_f=\theta_\pi$ on the elliptic set
(elliptic isolation). Then $0=J(f)-\mgeom(\theta_f)=\mathrm{defect}[\pi]$. Finiteness is multiplicity one.
\end{proof}

\begin{theorem}[The geometric expansion under (HYP)]
\label{geomAt_of_hyp}
\lean{GSpinGGP.geomAt_of_hyp}
\leanok
\uses{CentralCharOn, GGPData, GeomAt, HYP}
(Thesis Thm.~12.2.1 under (HYP).) If (HYP) holds for $G$, then $J(f)=\mgeom(\theta_f)$ for every central
character $\chi$ and every $f\in\scusp(A_G\backslash G,\chi^{-1})$.
\end{theorem}
\begin{proof}
\leanok
\uses{G, GSpin, GammaClass, H, IsGammaElt, IsQuasiCharacter, LIT-germ_unique, LIT-hc_regularSet_open_conull, N, QC, ScuspTestFn, deltaGGP, distribJ, distribJ_add, geomAt_away, geomAt_near, germ, gspin, isThetaF_chiProj, jHat, mGeom, mGeom_add, mem_mul_parabolicU_iff, orbitChoice, thetaF, thetaF_mem_QC, thetaF_spec, weylDisc}
Choose the invariant clopen $\Omega\ni1$ of the expansion near $1$; cut $f=f_1+(f-f_1)$ with $f_1$ supported in $\Omega$
and $\theta_{f_1}=1_\Omega\theta_f$; apply the expansion near $1$ to $f_1$ and the expansion away from $1$ to
$f-f_1$, and add, using linearity of $J$ and $\mgeom$ (Chapter~\ref{chap:geometric}).
\end{proof}

\begin{theorem}[The induction]
\label{geom_mult_all}
\lean{GSpinGGP.geom_mult_all}
\leanok
\uses{CentralCharOn, GGPData, GeomAt, MultAt}
For every GGP triple, the geometric expansion (for every central character) and the multiplicity formula hold.
\end{theorem}
\begin{proof}
\leanok
\uses{HYP, geomAt_of_hyp, multAt_of_hyp_geom}
Strong induction on $\dim G$: (HYP) is the induction hypothesis; apply the two preceding results.
\end{proof}

\begin{theorem}[Multiplicity formula from the trace formula]
\label{multiplicityFormula_of_traceFormula}
\lean{GSpinGGP.multiplicityFormula_of_traceFormula}
\leanok
\uses{MultiplicityFormula}
$m(\pi)=\mgeom(\theta_\pi)$ (and $m(\pi)<\infty$) for every GGP triple and every tempered $\pi$ trivial on
$A_{G,H}(F)$ (Thesis Thm.~7.2.2 = Thm.~1.0.2).
\end{theorem}
\begin{proof}
\leanok
\uses{CentralCharOn, GGPData, GeomAt, MultAt, geom_mult_all}
The multiplicity-formula half of the induction.
\end{proof}

